\section{C5 Surface Screen for C4 Projective Compaction}
\label{app:c5-surface}

C4 projective compaction reduces the full C2 dictionary from $q^b-1$ nonzero vectors per block to
\[
  \frac{q^b-1}{q-1}
\]
projective representatives.  This improves storage and reduces the public CLPN sample surface by a factor close to $q$.  However, a complete projective set still contains many linear dependencies.

\paragraph{Heuristic relation weight.}
For a full projective representative set in a block of width $b$, same-line weight-2 relations are removed.  A generic lower-bound screen therefore treats the first unavoidable relation weight as
\[
  b+1,
\]
because any $b+1$ vectors in a $b$-dimensional vector space are linearly dependent.  This is not a proof of security; it is a diagnostic for public-sample exposure.  The C5 estimator records this heuristic relation weight, the number of public CLPN rows, and the q-ary piling-up noise of relation combinations.

\paragraph{Relation-noise screen.}
If a relation combines $r$ CLPN ciphertexts whose plaintext masks sum to zero, then the combined ciphertext is an encryption of zero with q-ary error rate approximately
\[
  \eta_r=\frac{q-1}{q}\left(1-\left(1-\frac{q\eta_c}{q-1}\right)^r\right).
\]
A low $r$ and low $\eta_r$ create a useful warning surface for BKW-style, information-set-decoding, or large-sample q-ary-LPN analysis.  C5 flags width-2 projective blocks because their heuristic relation weight is $3$.

\paragraph{Compaction-noise screen.}
If a final row activates $L$ blocks and C4 uses at most one projective representative per active block, the compaction layer uses at most $L$ CLPN terms.  The C5 estimator reports
\[
  \eta_L=\frac{q-1}{q}\left(1-\left(1-\frac{q\eta_c}{q-1}\right)^L\right)
\]
and a conservative repetition-code failure bound.  This is a correctness diagnostic only; a production construction should replace repetition decoding by a stronger q-ary code.

\paragraph{C5 conclusion.}
C4 is better than C2/C3 with respect to dictionary size and public entries, but C5 shows that complete projective dictionaries still expose structured relation surfaces.  A future C6 direction should either increase block width, use partial randomized projective covers, add CRT-coded compaction to reduce active blocks, or replace the repetition-code CLPN layer by a stronger q-ary code with a full large-sample LPN analysis.
